% Запасные (дополнительные) слайды

\section{Приложение}

\begin{frame}{Архитектура системы обнаружения}
  \centering
  \begin{tikzpicture}[node distance=1.4cm and 2cm, >=Stealth, thick,
      block/.style={rectangle, draw=black, fill=gray!10, rounded corners, minimum width=2.2cm, minimum height=0.8cm, align=center}
  ]
  \node[block] (input) {Сырые логи\\(Sysmon, Zeek, EDR)};
  \node[block, right=of input] (parser) {Парсер и\\нормализация};
  \node[block, right=of parser] (graph) {Построение\\атрибутивного\\метаграфа};
  \node[block, right=of graph] (detect) {Сопоставление с\\тактиками ATT\&CK};
  \node[block, right=of detect] (alert) {Генерация\\инцидента};

  \draw[->] (input) -- (parser);
  \draw[->] (parser) -- (graph);
  \draw[->] (graph) -- (detect);
  \draw[->] (detect) -- (alert);
  \end{tikzpicture}
\end{frame}

\begin{frame}{Сравнение по времени обработки}
  \centering
  \begin{tabular}{lcc}
    \toprule
    Метод & Время (1 млн событий) & Память \\
    \midrule
    Snort & 3.1 с & 0.4 ГБ \\
    Isolation Forest & 12.7 с & 2.1 ГБ \\
    GNN & 24.5 с & 4.8 ГБ \\
    Предложенный & \textbf{8.2 с} & \textbf{1.1 ГБ} \\
    \bottomrule
  \end{tabular}
  \vspace{0.5cm}
  \footnotesize Тестирование: Intel i7-12700, 32 ГБ RAM, Ubuntu 22.04
\end{frame}

\begin{frame}{Поддерживаемые тактики MITRE ATT\&CK}
    \begin{itemize}
    \item TA0001 — Initial Access
    \item TA0002 — Execution
    \item TA0003 — Persistence
    \item TA0004 — Privilege Escalation
    \item TA0005 — Defense Evasion
    \item TA0007 — Discovery
    \item TA0008 — Lateral Movement
    \item TA0010 — Exfiltration
    \item TA0011 — Command and Control
    \end{itemize}
\end{frame}

\begin{comment}
\begin{frame}
    \frametitle{Ответы на замечания ведущей организации НИИ~<<Рога~и~копыта>>}
    \begin{itemize}
        \item Замечание -- ответ
        \item Замечание -- ответ
        \item Замечание -- ответ
        \item Замечание -- ответ
        \item Замечание -- ответ
    \end{itemize}
\end{frame}

\begin{frame}
    \frametitle{Ответы на замечания оф. оппонента Иванова\,И.\,И}
    \begin{itemize}
        \item Замечание -- ответ
        \item Замечание -- ответ
        \item Замечание -- ответ
        \item Замечание -- ответ
        \item Замечание -- ответ
    \end{itemize}
\end{frame}

\begin{frame}
    \frametitle{Ответы на замечания Петрова\,П.\,П}
    \begin{itemize}
        \item Замечание -- ответ
        \item Замечание -- ответ
        \item Замечание -- ответ
        \item Замечание -- ответ
        \item Замечание -- ответ
    \end{itemize}
\end{frame}
\end{comment}